% =====================================================================
\chapter{Momentum}
\label{tema:4}
% =====================================================================

El \cref{tema:2} mostró que el descenso de gradiente es lento en valles alargados: la dirección de curvatura alta obliga a usar pasos pequeños y la de curvatura baja avanza despacio. El método de momentum, propuesto por Polyak~\cite{polyak} con el nombre de método de la bola pesada, corrige este defecto acumulando los gradientes pasados. Este tema lo presenta como una media móvil exponencial, lo interpreta como la discretización de un sistema mecánico con rozamiento y demuestra que, en cuadráticas, reduce el número de iteraciones de un orden $\kappa$ a un orden $\sqrt\kappa$. Se sigue~\cite[§12.6]{d2l}, \cite[§8.3.2 y §8.3.3]{gbc}, la clase~23 de~\cite{mit18065} y el artículo divulgativo de Goh~\cite{goh}.

\section{Medias móviles exponenciales}

\begin{definicion}[Media móvil exponencial]\label{def:ema}
Sean $g_1,g_2,\dots$ vectores y $\beta\in[0,1)$. La \emph{acumulación exponencial} con $v_0=0$ es $v_k=\beta v_{k-1}+g_k$, y la \emph{media móvil exponencial} es $\bar v_k=\beta\bar v_{k-1}+(1-\beta)g_k$, con $\bar v_0=0$.
\end{definicion}

\begin{proposicion}\label{prop:ema}
Con la notación de la \cref{def:ema},
\[
v_k=\sum_{j=1}^{k}\beta^{k-j}g_j,\qquad \bar v_k=(1-\beta)\,v_k.
\]
En particular, si $g_j=g$ para todo $j$, entonces $v_k=\frac{1-\beta^{k}}{1-\beta}\,g\to\frac{g}{1-\beta}$ y $\bar v_k=(1-\beta^{k})\,g\to g$.
\end{proposicion}

\begin{proof}
La fórmula de $v_k$ se demuestra por inducción: es cierta para $k=1$, y si lo es para $k-1$, entonces $v_k=\beta\sum_{j=1}^{k-1}\beta^{k-1-j}g_j+g_k=\sum_{j=1}^{k}\beta^{k-j}g_j$. Las dos recurrencias son lineales con la misma condición inicial y se diferencian en el factor $(1-\beta)$ del término nuevo, luego $\bar v_k=(1-\beta)v_k$. Con $g_j=g$, la suma es geométrica: $\sum_{j=1}^{k}\beta^{k-j}=\sum_{i=0}^{k-1}\beta^{i}=(1-\beta^{k})/(1-\beta)$.
\end{proof}

El peso de un gradiente decae geométricamente con su antigüedad; la suma total de los pesos es $1/(1-\beta)$, que se interpreta como el número efectivo de gradientes promediados~\cite[§12.6.1]{d2l}. Con $\beta=0{,}9$ son unos~10; con $\beta=0{,}99$, unos~100.

\section{El método de la bola pesada}

\begin{definicion}[Momentum]\label{def:momentum}
Dados $\theta_0$, $v_0=0$, un paso $\eta>0$ y un coeficiente $\beta\in[0,1)$, el método de \emph{momentum} (o de la \emph{bola pesada}) es
\[
v_{k+1}=\beta v_k+\nabla f(\theta_k),\qquad \theta_{k+1}=\theta_k-\eta\,v_{k+1}.
\]
Con $\beta=0$ se recupera el descenso de gradiente.
\end{definicion}

\begin{proposicion}[Forma de dos pasos]\label{prop:dos-pasos}
Los iterados del método de momentum cumplen, para $k\ge 1$,
\[
\theta_{k+1}=\theta_k-\eta\nabla f(\theta_k)+\beta\,(\theta_k-\theta_{k-1}).
\]
\end{proposicion}

\begin{proof}
De la segunda ecuación, $\eta v_k=\theta_{k-1}-\theta_k$. Sustituyendo la primera en la segunda, $\theta_{k+1}=\theta_k-\eta\beta v_k-\eta\nabla f(\theta_k)=\theta_k-\eta\nabla f(\theta_k)+\beta(\theta_k-\theta_{k-1})$.
\end{proof}

El término $\beta(\theta_k-\theta_{k-1})$ es la inercia: el iterado conserva una fracción $\beta$ de su desplazamiento anterior. Si el gradiente es constante, la \cref{prop:ema} da $\eta v_k\to\eta g/(1-\beta)$: con $\beta=0{,}9$, el paso efectivo es diez veces el del descenso de gradiente. En un valle estrecho, las componentes del gradiente que cambian de signo de un paso al siguiente (las transversales al valle) se cancelan parcialmente en $v_k$, mientras que las que conservan el signo (a lo largo del valle) se acumulan.

\begin{ejemplo}\label{ej:momentum}
Sea $f(x)=x^{2}/2$, con $f'(x)=x$, $x_0=1$, $\eta=0{,}1$ y $\beta=0{,}9$:
\begin{center}\small
\begin{tabular}{@{}ccccc@{}}
\toprule
$k$ & $v_k$ & $x_k$ (momentum) & $x_k$ (descenso de gradiente) \\
\midrule
1 & $1$ & $0{,}9$ & $0{,}9$ \\
2 & $1{,}8$ & $0{,}72$ & $0{,}81$ \\
3 & $2{,}34$ & $0{,}486$ & $0{,}729$ \\
4 & $2{,}592$ & $0{,}2268$ & $0{,}6561$ \\
5 & $2{,}5596$ & $-0{,}02916$ & $0{,}59049$ \\
\bottomrule
\end{tabular}
\end{center}
El momentum avanza más deprisa y en el paso~5 sobrepasa el mínimo: la bola tiene inercia.
\end{ejemplo}

\section{Interpretación física}

Considérese una partícula de masa $m>0$ que se mueve en un potencial $f$ con rozamiento viscoso de coeficiente $\gamma>0$. Su posición $\theta(t)$ y su velocidad $u(t)=\dot\theta(t)$ cumplen
\begin{equation}\label{eq:bola}
m\,\dot u=-\nabla f(\theta)-\gamma\,u,\qquad\dot\theta=u.
\end{equation}

\begin{proposicion}[Disipación de energía]\label{prop:energia}
Sea $E(t)=\frac{m}{2}\norm{u(t)}^{2}+f(\theta(t))$ la energía mecánica. A lo largo de las soluciones de~\eqref{eq:bola}, $\dot E=-\gamma\norm{u}^{2}\le 0$.
\end{proposicion}

\begin{proof}
Por la regla de la cadena, $\dot E=m\,u\T\dot u+\nabla f(\theta)\T\dot\theta=u\T\bigl(-\nabla f(\theta)-\gamma u\bigr)+\nabla f(\theta)\T u=-\gamma\norm{u}^{2}$.
\end{proof}

La energía decrece mientras la partícula se mueve y solo se estabiliza cuando $u=0$, es decir, en reposo en un punto crítico. Sin rozamiento, la energía se conserva y la partícula oscila indefinidamente.

\begin{proposicion}[Discretización]\label{prop:discretizacion}
Se discretiza~\eqref{eq:bola} con paso temporal $\Delta t>0$ mediante
\[
u_{k+1}=u_k+\frac{\Delta t}{m}\bigl(-\nabla f(\theta_k)-\gamma u_k\bigr),\qquad\theta_{k+1}=\theta_k+\Delta t\,u_{k+1}.
\]
Con $v_k=-(m/\Delta t)\,u_k$, la iteración es exactamente el método de momentum de la \cref{def:momentum} con
\[
\beta=1-\frac{\gamma\,\Delta t}{m},\qquad\eta=\frac{\Delta t^{2}}{m}.
\]
\end{proposicion}

\begin{proof}
Multiplicando la primera ecuación por $-m/\Delta t$: $v_{k+1}=(1-\gamma\Delta t/m)\,v_k+\nabla f(\theta_k)=\beta v_k+\nabla f(\theta_k)$. Y $\Delta t\,u_{k+1}=-(\Delta t^{2}/m)\,v_{k+1}=-\eta v_{k+1}$.
\end{proof}

El coeficiente $\beta$ mide la fracción de velocidad que sobrevive a un paso: $\beta\to 1$ corresponde a rozamiento nulo y $\beta=0$ a un rozamiento tan grande que la partícula olvida su velocidad en cada paso, que es el descenso de gradiente.

\section{Convergencia en cuadráticas}

\begin{lema}[Una dirección propia]\label{lem:polinomio}
Sea $f(z)=\frac{\lambda}{2}z^{2}$ con $\lambda>0$. Los iterados del método de momentum cumplen $z_{k+1}=(1+\beta-\eta\lambda)\,z_k-\beta z_{k-1}$, cuyo polinomio característico es
\[
p(r)=r^{2}-(1+\beta-\eta\lambda)\,r+\beta.
\]
Si $(1-\sqrt\beta)^{2}\le\eta\lambda\le(1+\sqrt\beta)^{2}$, las dos raíces de $p$ tienen módulo $\sqrt\beta$.
\end{lema}

\begin{proof}
La recurrencia es la \cref{prop:dos-pasos} con $\nabla f(z)=\lambda z$. Sea $b=1+\beta-\eta\lambda$. El discriminante de $p$ es $b^{2}-4\beta$, que es $\le 0$ si y solo si $\abs{b}\le 2\sqrt\beta$, es decir, $1+\beta-2\sqrt\beta\le\eta\lambda\le 1+\beta+2\sqrt\beta$, que es la condición del enunciado. En ese caso las raíces son complejas conjugadas (o una raíz doble real), y su producto, igual a $\abs{r}^{2}$, es el término independiente $\beta$.
\end{proof}

\begin{teorema}[Momentum óptimo en cuadráticas]\label{teo:heavy-ball}
Sea $f(\theta)=\frac12\theta\T Q\theta$ con $Q$ simétrica de autovalores en $[\mu,L]$, $0<\mu<L$, y $\kappa=L/\mu$. Con
\[
\eta=\frac{4}{(\sqrt L+\sqrt\mu)^{2}},\qquad\beta=\Bigl(\frac{\sqrt\kappa-1}{\sqrt\kappa+1}\Bigr)^{2},
\]
la iteración lineal del método de momentum sobre el par $(\theta_k,\theta_{k-1})$ tiene radio espectral
\[
\rho=\sqrt\beta=\frac{\sqrt\kappa-1}{\sqrt\kappa+1}.
\]
\end{teorema}

\begin{proof}
En la base de autovectores de $Q$ la iteración se desacopla, como en el \cref{teo:gd-cuadratica}, en $n$ recurrencias escalares como la del \cref{lem:polinomio}, una por autovalor $\lambda_i$. Cada una es una iteración lineal en $\R^{2}$ cuya matriz tiene por autovalores las raíces de $p$; el radio espectral global es el máximo de los módulos de esas raíces. Se comprueba que todos los $\lambda\in[\mu,L]$ cumplen la condición del lema. Como $\sqrt\beta=(\sqrt L-\sqrt\mu)/(\sqrt L+\sqrt\mu)$,
\[
1-\sqrt\beta=\frac{2\sqrt\mu}{\sqrt L+\sqrt\mu},\qquad 1+\sqrt\beta=\frac{2\sqrt L}{\sqrt L+\sqrt\mu},
\]
de donde $(1-\sqrt\beta)^{2}=4\mu/(\sqrt L+\sqrt\mu)^{2}=\eta\mu$ y $(1+\sqrt\beta)^{2}=\eta L$. Para $\lambda\in[\mu,L]$, $\eta\lambda\in[\eta\mu,\eta L]$, que es justamente el intervalo del lema. Todas las raíces tienen módulo $\sqrt\beta$.
\end{proof}

\begin{observacion}
El radio espectral determina la tasa asintótica: para todo $\delta>0$ existe $C$ tal que $\norm{\theta_k}\le C(\rho+\delta)^{k}$. A diferencia del \cref{teo:gd-cuadratica}, la matriz de la iteración no es simétrica y la constante $C$ puede ser mayor que~1; en los extremos $\lambda=\mu$ y $\lambda=L$ la raíz es doble y aparece un factor polinómico en $k$. Por eso las iteraciones reales pueden ser algo más que las que predice $\rho$ (\cref{eje:4-autovalores}).
\end{observacion}

Comparando con el \cref{teo:gd-cuadratica}, para $\kappa$ grande, $\ln\rho_{\text{mom}}\approx -2/\sqrt\kappa$ frente a $\ln\rho_{\text{gd}}\approx -2/\kappa$: el número de iteraciones pasa de ser proporcional a $\kappa$ a serlo a $\sqrt\kappa$. Con $\kappa=100$, el factor baja de $0{,}9802$ a $0{,}8182$ y las iteraciones para reducir el error a $10^{-3}$, de 346 a 35 (\cref{tab:kappa}). El artículo de Goh~\cite{goh} permite explorar este resultado de forma interactiva.

\begin{ejemplo}[Rosenbrock]
En la función del \cref{ej:rosenbrock}, desde $(-1{,}5;\,2)$ y con $\eta=10^{-3}$, el método de momentum con $\beta=0{,}9$ alcanza $f<10^{-4}$ en la iteración~934, mientras que el descenso de gradiente con el mismo paso no lo logra en 5000 iteraciones.
\end{ejemplo}

\section{Momentum de Nesterov}

\begin{definicion}[Momentum de Nesterov]
El método de \emph{momentum de Nesterov} evalúa el gradiente en el punto al que llevaría la inercia:
\[
v_{k+1}=\beta v_k+\nabla f(\theta_k-\eta\beta v_k),\qquad\theta_{k+1}=\theta_k-\eta v_{k+1}.
\]
\end{definicion}

Para funciones convexas y $L$-suaves, una versión de este método con $\beta$ variable alcanza $f(\theta_k)-f^{\ast}=O(1/k^{2})$, frente al $O(1/k)$ del descenso de gradiente, y ese orden es óptimo entre los métodos de primer orden~\cite{nesterov}. En el caso estocástico, en cambio, el momentum de Nesterov no mejora el orden de convergencia~\cite[§8.3.3]{gbc}. En el entrenamiento de redes profundas, el momentum se usa sobre todo a través de Adam (\cref{tema:5}), cuyo primer momento es una media móvil exponencial del gradiente.

\begin{lecturas}
\item Zhang \emph{et al.}~\cite{d2l}, §12.6 «Momentum», en particular §12.6.3 «Theoretical Analysis».
\item Goodfellow, Bengio y Courville~\cite{gbc}, §8.3.2 «Momentum» y §8.3.3 «Nesterov Momentum».
\item MIT 18.065~\cite{mit18065}, clase~23 «Accelerating Gradient Descent (Use Momentum)».
\item Goh~\cite{goh}, «Why Momentum Really Works», exposición interactiva del análisis de este tema.
\item Polyak~\cite{polyak}, artículo original del método de la bola pesada.
\end{lecturas}

\section*{Ejercicios del tema 4}
\addcontentsline{toc}{section}{Ejercicios del tema 4}

\begin{ejercicio}{1}{eje:4-efectivo}
Con gradiente constante $g$, ¿por qué factor multiplica el momentum el paso del descenso de gradiente a largo plazo si $\beta=0{,}9$? ¿Y si $\beta=0{,}99$? ¿Cuántos pasos hacen falta, con $\beta=0{,}9$, para alcanzar el 90~\% de ese valor? Explique por qué un $\beta$ grande favorece que el iterado sobrepase el mínimo.
\end{ejercicio}

\begin{ejercicio}{2}{eje:4-pasos}
Aplique tres pasos del método de momentum a $f(x)=x^{2}/2$ con $x_0=1$, $v_0=0$, $\eta=0{,}1$ y $\beta=0{,}9$, y compárelos con tres pasos de descenso de gradiente con el mismo $\eta$. ¿Qué fracción de la distancia inicial al mínimo ha recorrido cada método?
\end{ejercicio}

\begin{ejercicio}{2}{eje:4-optimos}
Una cuadrática tiene $\mu=1$ y $L=100$. Calcule los parámetros $\eta$ y $\beta$ del \cref{teo:heavy-ball}, el factor $\rho$ y el número de iteraciones estimado para reducir el error a $10^{-3}$. Compárelo con el descenso de gradiente con paso óptimo.
\end{ejercicio}

\begin{ejercicio}{2}{eje:4-autovalores}
Sea $f(x)=\frac12\sum_{i=1}^{3}\lambda_i x_i^{2}$ con $\lambda_i=2^{-(i-1)}$, es decir, $\lambda=(1,\tfrac12,\tfrac14)$, y $x_0=(1,1,1)$. (a) Con descenso de gradiente y $\eta=1$, ¿cuántos pasos hacen falta para que todas las coordenadas tengan módulo menor que $10^{-2}$? (b) Repita con el paso óptimo del \cref{teo:gd-cuadratica}. (c) Obtenga los parámetros óptimos del momentum, calcule a mano los dos primeros iterados y estime el número de pasos con el radio espectral.
\fuente{adaptado de Zhang \emph{et al.}~\cite{d2l}, §12.6, ejercicio~2 (autovalores $\lambda_i=2^{-i}$; aquí se reduce a tres coordenadas para el cálculo a mano).}
\end{ejercicio}

\begin{ejercicio}{3}{eje:4-forma}
Demuestre la \cref{prop:dos-pasos}. Demuestre después que, si $\nabla f(\theta_k)=g$ para todo $k$, el desplazamiento cumple $\theta_k-\theta_{k+1}=\eta\,\frac{1-\beta^{k+1}}{1-\beta}\,g$.
\end{ejercicio}

\begin{ejercicio}{3}{eje:4-estabilidad}
Sea $f(z)=\frac{\lambda}{2}z^{2}$ con $\lambda>0$, $\eta>0$ y $\beta\in[0,1)$. Demuestre que las dos raíces del polinomio $p$ del \cref{lem:polinomio} tienen módulo menor que~1 si y solo si $0<\eta\lambda<2(1+\beta)$. Compare con el límite $\eta\lambda<2$ del descenso de gradiente.
\end{ejercicio}
